\section{框架与方法}
\reviewslide[先确定页面要表达的判断，再选择图文形式。]{从材料到页面形成连续表达}{材料经过问题提炼和证据筛选，转化为可以口头讲解的页面。}{
  \centering
  \begin{tikzpicture}[
    box/.style={draw=TJBluePrimary!45,fill=TJBluePrimary!5,rounded corners=2pt,
      minimum width=3.5cm,minimum height=1.35cm,align=center,font=\small},
    >=Latex]
    \node[box] (a) {原始材料\\论文、报告与数据};
    \node[box,right=1cm of a] (b) {问题与证据\\保留主张及适用条件};
    \node[box,right=1cm of b] (c) {汇报页面\\主旨、主体与出处};
    \draw[->,draw=TJBluePrimary,line width=1pt] (a)--(b);
    \draw[->,draw=TJBluePrimary,line width=1pt] (b)--(c);
  \end{tikzpicture}
}{自绘流程示意；用于验证图文页与底部结论条。}

\section{结果与比较}
\reviewslide{不同内容使用对应的主体形式}{比较内容采用表格，保留相同的判断维度。}{
  \small\renewcommand{\arraystretch}{1.65}
  \begin{tabularx}{\linewidth}{@{}lXX@{}}
    \toprule
    内容类型 & 主体形式 & 阅读重点 \\
    \midrule
    机制与流程 & 框架图或流程图 & 环节之间的联系 \\
    结果与对照 & 图表或比较表 & 条件、指标与差异 \\
    定义与推导 & 公式与符号说明 & 含义、假设与边界 \\
    \bottomrule
  \end{tabularx}
}{版式示例：内容类型、表达形式与阅读重点。}
